\subsection{Separable extensions}

Let K be a field. Since any extension of K is a K-algebra, the notion of separability introduced in Def. 1 (V, p.~119) applies in particular to the case of extensions of K. By example 3 (V, p.~119) that definition of separability agrees in the case of algebraic extensions with that of § 7 (V, p.~36, Def. 1).

PROPOSITION 6. --- Every pure extension of a field K is separable.

This follows at once from Example 1 (V, p.~119) and Prop. 4 (V, p.~120).

PROPOSITION 7. --- Let E be a field, G a group of automorphisms of E and K the subfield of E consisting of the invariants of G. Then E is a separable extension of K.

Let L be an extension of K; there exists an algebraically closed extension \(\Omega\) of L whose transcendence degree over K is at least equal to that of E over K. By Cor. 1 (V, p.~117) there exists a K-homomorphism \(u\) of E into \(\Omega\) . We denote by \(v\) the \(\Omega\) -algebra homomorphism of \(\mathbf{A} = \Omega \otimes_{\mathbf{K}} \mathbf{E}\) into \(\Omega\) which maps \(\lambda \otimes x\) to \(\lambda \cdot u(x)\) for \(\lambda \in \Omega\) and \(x \in \Omega\) ; we write \(\mathfrak{a}\) for the kernel of \(v\) .

For every \(s \in G\) let \(h_s\) be the automorphism \(\operatorname{Id}_{\Omega} \otimes s\) of the \(\Omega\) -algebra \(A\) ; the kernel of the homomorphism \(v \circ h_s\) of \(A\) into \(\Omega\) is the prime ideal \(\mathfrak{a}_s = h_s^{-1}(\mathfrak{a})\) of \(A\) . It is clear that the ideal \(\mathfrak{b} = \bigcap_{s \in G} \mathfrak{a}_s\) of \(A\) is stable under the automorphisms

\(h_{s}\) . Therefore (V, p.~63, Cor.) the ideal b is of the form \(c \otimes_{K} E\) , where c is an ideal of \(\Omega\) . Now \(b \subset a \neq A\) , and so \(c \neq \Omega\) ; since \(\Omega\) is a field, we have c = 0, hence b = 0.

The family \((\mathfrak{a}_{s})_{s\in G}\) of prime ideals of A thus has zero intersection. By Prop. 2 (V, p.~118) the ring is reduced, and a fortiori the same holds of the subring \(L\otimes_{K}E\) of A. Since L was an arbitrary extension of K, this proves that E is separable over K.

PROPOSITION 8. --- Let L be an extension of a field K. If L is separable over K, every subextension of L is separable over K. Conversely, if every finitely generated subextension of L is separable over K, then L is separable over K.

This follows at once from Prop. 3, a) and b) (V, p.~119).

Thus separability may be called a property of « finite character ».

PROPOSITION 9. --- Let L be an extension of a field K and M a commutative L-algebra (for example, an extension of L). If M is separable over L and L is separable over K, then M is separable over K.

Let \(K'\) be an extension of K. Since L is a separable extension of K, the ring \(K' \otimes_K L\) is reduced; since M is a separable L-algebra, Prop. 5 (V, p.~120) shows the ring \((K' \otimes_K L) \otimes_L M\) to be reduced. Now the ring \(K' \otimes_K M\) is isomorphic to \((K' \otimes_K L) \otimes_L M\) (II, p.~278, Prop. 2), hence is reduced. This proves that M is separable over K.

\pandocbounded{\includegraphics[keepaspectratio]{images/part002_8e27f267.jpg}}

If the extension M is separable over K, it is not necessarily separable over L (cf.~however V, p.~124, Cor. 3). For example, if p is a prime number, the field \(\mathbf{F}_{p}(\mathbf{X})\) of rational fractions in one indeterminate X over \(F_{p}\) is separable over \(F_{p}\) (V, p.~121, Prop. 6) but it is a p-radical algebraic extension of \(\mathbf{F}_{p}(\mathbf{X}^{p})\) ; in particular \(\mathbf{F}_{p}(\mathbf{X})\) is not separable over \(\mathbf{F}_{p}(\mathbf{X}^{p})\) .

Later (V, p.~137 Prop. 5) we shall study the separability of composite extensions.
% semantic-fragments: legacy-input-seam
\relax{}
